\documentclass[hyperref={pdfpagelabels=false},table]{beamer}

\input{lec_common}
\date{}
\begin{document}

\part{Unit 7: Exceptions, File Handling, Iterable objects}

%\frame{\tableofcontents}


\begin{frame}[fragile]{Exceptions}
\begin{block}{Creating errors}
Creating an error is called ``raising an exception''. You may raise an exception like this:
\begin{python}
raise Exception("Something went wrong")
\end{python}
\end{block}
Typical exceptions are 
\begin{itemize}
	\item \pyth!TypeError!
	\item \pyth!ValueError!
	\item \pyth!IndexError!
\end{itemize}

You already know \pyth{SyntaxError} and \pyth{ZeroDivisionError}.
\end{frame}


\begin{frame}[fragile]{Review the alarms}
\begin{block}{Reviewing the errors}
You may review the errors using \pyth{try} and \pyth{except}:
\begin{python}
try:
  <some code that might raise an exception>
  # like raise Exception('Something went wrong')
  # other exceptions possible 
except ValueError:
  print('Oops, a ValueError occurred')
except TypeError:
  print('We got a TypeError but not a ValueError')
except Exception:
  print('Some other kind of exception occurred.')
finally: 
  print('Everything good!')
\end{python}
\end{block}

\begin{alertblock}{Flow Control}
An exception stops the flow and looks for the closest enclosing \pyth{try} block. If it is not caught it continues searching for the next \pyth{try} block.
\end{alertblock}

\end{frame}


%%%%%%%%%%%%%%%%%%%%%%%%%%%%%%%%%%%%%%%%%%%%%%%%%%%%%%%%%%%%%%%%%%%%%
%%%%%%%%%%%%%%%%%%%%%%%%%%%%%%%%%%%%%%%%%%%%%%%%%%%%%%%%%%%%%%%%%%%%%
\begin{frame}\frametitle{Inheritance and class hierarchies}

Example: Hierarchy of some built-in exceptions classes:
\vskip 10pt 
  \includegraphics[width=1.05\textwidth]{exceptions}

\end{frame}

\begin{frame}[fragile]{Error messages}
\begin{alertblock}{Golden rule}
Never print error messages, \alert{raise an exception instead}!
\end{alertblock}
\vfill

Don't do this:
\begin{python}
def fixpointIter(f, x0, maxit = 100, tol = 1e-6):
  x = copy.copy(x0)
  for i in range(maxit):
    fx = f(x)
    if abs(fx-x) < tol:
      break
    x = fx
  else:
    print(
     f"No convergence in {maxit} iterations.")
  return x, i
\end{python}
\end{frame}
\begin{frame}[fragile]{Error messages (Cont.)}
Do this instead:
\begin{python}
def fixpointIter(f, x0, maxit = 100, tol = 1e-6):
  x = copy.copy(x0)
  for i in range(maxit):
    fx = f(x)
    if abs(fx-x) < tol:
      break
    x = fx
  else:
    raise Exception(
        f"No convergence in {maxit} iterations.")
  return x, i
\end{python}
\end{frame}

\begin{frame}[fragile]{Catching the error}
With the last construction we can do things like this:
\begin{python}
for it in [10, 100, 1000, 10000]:
  try:
    x, numit = fixpointIter(f, x0, maxit = it)
  except Exception:
    print(f"{it} iterations were not enough.")
  else:
    print(f"{it} iterations were enough!")
    print(f"Converged to {x} in {numit} iterations")
    break    
\end{python}


\end{frame}

\begin{frame}[fragile]{Define your own Exceptions}
You can construct your own expeption by using simple inheritance:
\vfill
\begin{python}
 class NoConvergenceException(Exception):
    pass
\end{python}
\vfill
Everything but the name is inherited her from the parent class.

\end{frame}

%%%%%%%%%%%%%%%%%%%%%%%%%%%%%%%%%%%%%%%%%%%%%%%%%%%%%%%%%%%%%%%%%%%%%
%%%%%%%%%%%%%%%%%%%%%%%%%%%%%%%%%%%%%%%%%%%%%%%%%%%%%%%%%%%%%%%%%%%%%
\begin{frame}[fragile]\frametitle{split and join}

  The methods \pyth{split} and \pyth{join} can be used on instances of \pyth{str}.

\begin{python}
separation.join(some_list)    # returns a string
\end{python}

\pause

\begin{python}
some_string.split(separation) # returns a list
\end{python}

In both cases, separation is a string.

\end{frame}

%%%%%%%%%%%%%%%%%%%%%%%%%%%%%%%%%%%%%%%%%%%%%%%%%%%%%%%%%%%%%%%%%%%%%
%%%%%%%%%%%%%%%%%%%%%%%%%%%%%%%%%%%%%%%%%%%%%%%%%%%%%%%%%%%%%%%%%%%%%
\begin{frame}[fragile]\frametitle{Examples using join}

The methods split and join can be used on instances of str.

\begin{python}
town_list = ["Stockholm", "Paris", "London"]
towns1 = "".join(town_list)
towns2 = "\t".join(town_list)  # tab
towns3 = ",".join(town_list) 
print(towns)  # StockholmParisLondon
print(towns2)  # Stockholm	Paris	London
print(towns3)  # Stockholm,Paris,London
\end{python}

\pause
When the list elements are not strings, they must be cast to \pyth{str}.

\begin{python}
nr_list = [3, 7.4, 1e-4, 3+7j]
numbers = " and ".join(str(nr) for nr in nr_list)
print(numbers) # 3 and 7.4 and 0.0001 and (3+7j)   
\end{python}

  In both cases, separation is a string (\pyth{str}).
\end{frame}

%%%%%%%%%%%%%%%%%%%%%%%%%%%%%%%%%%%%%%%%%%%%%%%%%%%%%%%%%%%%%%%%%%%%%
%%%%%%%%%%%%%%%%%%%%%%%%%%%%%%%%%%%%%%%%%%%%%%%%%%%%%%%%%%%%%%%%%%%%%
\begin{frame}[fragile]\frametitle{Examples using split}

\begin{python}
msg ="Let no one ignorant of geometry enter!"
L1 = msg.split()
print(L1)
# ["Let", "no", "one", "ignorant", "of", "geometry", "enter!"]

L2 = msg.split("of geometry")
print(L2) # ["Let no one ignorant ", " enter!"] 
\end{python}

\begin{python}
filename = "test.min.js"
L3 = filename.split(".")
print("The file extension is", L3[-1]) 
# The file extension is js
\end{python}

\end{frame}


\begin{frame}[fragile]{File I/O}
 File I/O (in- and output) is essential when 
\begin{itemize}
\item working with measured or scanned data
\item interacting with other programs 
\item saving information for comparisons or other postprocessing needs
\item .....
\end{itemize}
\end{frame}
\begin{frame}[fragile]{File objects}
A file is a Python object with associated methods:
\begin{python}
 # Creating a read-only file object
 myfile=open("measurement.dat","r") 
\end{python}

The whole file can be read and stored in a string by
\begin{python}
  s = myfile.read()
\end{python}
You can also read it one line at a time by
\begin{python}
myfile.readline()
\end{python}
or like
\begin{python}
for line in myfile.readlines()
  print(line)
\end{python}
\end{frame}


\begin{frame}[fragile]{Files as generators}
A file object is a \alert{generator}. We will talk more about generators later.
\vfill
A generator is like a list, except the values need not exist until asked
for.
\vfill
A main feature of generators is that they are \emph{disposable}.
When you read a line from a file, it is removed from the file object (not from the file itself).
The following code will print three different things:

\begin{python}
print(myfile.readline())  # Line 1
print(myfile.readline())  # Line 2
print(myfile.readline())  # Line 3
\end{python}
\end{frame}


\begin{frame}[fragile]{File close method}
 A file has to be closed before it can be reread.
\begin{python}
 myfile.close() # closes the file object
\end{python}
\vfill
It is automatically closed when 
\begin{itemize}
 \item the program ends
\item the enclosing program unit (e.g. function) is left .
\end{itemize}
\vfill
Before a file is closed, you won't see any changes in it
by an external editor!
\end{frame}

\begin{frame}[fragile]{The \pyth!with! statement}
If you forget to close a file, problems can occur. 
Also, an error might prevent you from closing the file. Consider
\begin{python}
myfile = open(name, "w")
myfile.write("some data")
a = 1/0
myfile.write("other data")
myfile.close()
\end{python}
The \pyth!with! statement helps with this:
\begin{python}
with open(name, "w") as myfile:
  myfile.write("some data")
  a = 1/0
  myfile.write("other data")
\end{python}
With this construction, the file is always closed even if an exception occurs. It is shorthand for a clever \pyth!try!-\pyth!except! block.
\end{frame}

\begin{frame}[fragile]{File Modes (read, write, etc.)}
 \begin{python}
   file1=open("file1.dat","r")  # read only
   file2=open("file2.dat","r+") # read/write
   file3=open("file3.dat","a")  # append
   file4=open("file4.dat","w")  # (over-)write
 \end{python}

The modes "r","r+","a" require that the file exists.
\begin{exampleblock}{File append example}
\begin{python}
  file3=open("file3.dat","a")
  file3.write("something new\n") # Note the "\n"
\end{python}
\end{exampleblock}
\end{frame}




\begin{frame}[fragile]{Saving NumPy arrays}
The read and write methods convert data to strings.

Complex data types (like arrays) cannot be written this way.

NumPy provides its own methods for storing arrays.
\begin{python}
a = array([10, 20, 30])
with open("outfile.txt", "w") as myfile:
  numpy.savetxt(myfile, a)  # Saves a in outfile.txt

with open("outfile", "wb") as myfile: # note "wb" for write binary
  numpy.save(myfile, a)  # Saves a in outfile (binary)
\end{python}

You can also just give the name of the file as a string:
\begin{python}
numpy.savetxt("outfile.txt", a)
# Saves a in outfile.txt
numpy.save("outfile", a) 
# Saves a in outfile.npy (binary)
\end{python}
  By using
  \href{https://numpy.org/doc/stable/reference/generated/numpy.savez.html}{numpy.savez}
  one can save several arrays in one file.
\end{frame}


\begin{frame}[fragile]{Reading NumPy arrays}
Reading arrays is done similarly
\begin{python}
with open("infile.txt", "r") as myfile:
  a = numpy.loadtxt(myfile)

with open("infile.npy", "r") as myfile:
  numpy.load(myfile, a)
\end{python}
or
\begin{python}
a = numpy.loadtxt("infile.txt")
a = numpy.load("infile.npy") 
\end{python}
If the data is separated by something other than spaces you should use the \pyth!delimiter! keyword
\begin{python}
a = numpy.loadtxt("infile.txt", delimiter = ";")
# Turns "1; 2; 3;" into array([1,2,3])
\end{python}
There are several more options like this, see the documentation.
\end{frame}
\begin{frame}[fragile]{Pickle, Shelve}
\vfill
Other methods to write and read data to a file are \pyth!pickle! and \pyth!shelve!.
We refer to the course book or other literature for those.
\vfill
\end{frame}
\begin{frame}[fragile]{Definition: Iterable objects}
\begin{defblock}
An iterable object has a special method \pyth!__iter__! which is called in for loops:
\end{defblock}

\begin{example}
A typical iterator is \pyth!range! in Python 3.x:
\begin{python}
for i in range(100000000):
   if i > 10:
    break
\end{python}

Is the same as 
\begin{python}
for i in range(100000000).__iter__():
  if i > 10:
    break
\end{python}
\pyth!range(10000).__iter__()! creates an \emph{iterator}.
\end{example}
\end{frame}

\begin{frame}[fragile]{Iterators }
\begin{defblock}
An iterator has a special method \pyth!__next__!.
\end{defblock}
\begin{example}
\begin{python}
rg=range(3)  # iterable object
rgi=rg.__iter__()  # an iterator
rgi.__next__()  # returns 0
rgi.__next__()  # returns 1
rgi.__next__()  # returns 2
rgi.__next__()  # returns StopIteration exception
\end{python}
\end{example}
Iterators can be exhausted. They exist as objects but are of no use any more
in your program.
\end{frame}
\begin{frame}[fragile]{Own iterators}
Creation of iterators is possible with the keyword \pyth!yield!:
\begin{python}
def odd_numbers(n):
  "generator for odd numbers less than n"
  for k in range(n):
    if k % 2 == 1:
      yield k
\end{python}

Then you call it as:
\begin{python}
g = odd_numbers(10)
for k in g:
  ... # do something with k
\end{python}

\end{frame}

\begin{frame}[fragile]{Iterator Tools}
\begin{itemize}
\item \pyth!enumerate! is used to \alert{enumerate} an iterable:
\begin{python}
A = ["a", "b", "c"]
for i, x in enumerate(A):
  print (i, x)  # result: 0 a 1 b 2 c
\end{python}
\item \pyth!reversed! creates an iterator from a list by going backwards:
\begin{python}
A = [0, 1, 2]
for elt in reversed(A):
  print(elt)    # result: 2 1 0
\end{python}
\end{itemize}
\end{frame}
\begin{frame}[fragile]{Infinite Iterators}
Iterators can be infinite:
\begin{example}
\begin{python}
def odd_numbers():
    "generator for odd numbers"
    k=0
    while True:
       k+=1
       if k % 2 == 1:
          yield k
      
# Find the first odd_number greater than 68:
on=odd_numbers()

for nu in on:
    if nu > 68:
       print(nu)
       break      
\end{python}
\end{example}
\end{frame}
\begin{frame}[fragile]{Example: Arithmetic Geometric Mean}
With $a_0 = 1$ and $b_0 = \sqrt{1-k^2}$, the iteration
\begin{align*}
  a_{i+1} &= \frac{a_i + b_i}{2} \\
  b_{i+1} &= \sqrt{a_ib_i}
\end{align*}
converges to
\begin{equation*}
  F(k, \pi/2) := \frac{\pi}{2} \lim_{i \to \infty}{\frac{1}{a_i}} = \int_0^{\frac{\pi}{2}}{\frac{1}{\sqrt{1-k^2\sin^2(\theta)}}\mathrm{d}\theta }
\end{equation*}
This is called a \emph{complete elliptic integral of the first kind}.
  \footnote{Legendre form of elliptic integrals of first, second and third
  kinds.}
\end{frame}
\begin{frame}[fragile]{Example: Arithmetic Geometric Mean (Cont.)}

\begin{python}
def arithmetic_geometric_mean(a,b):
  """
  Generator for the arithmetic and geometric mean
  a,b initial values
  """
  while True: # infinite loop
    a,b = (a+b)/2, sqrt(a*b)
    yield a,b
\end{python}

\begin{python}
def elliptic_integral(k, tolerance = 1e-5):
  """
  Compute an elliptic integral of the first kind.
  """
  a_0, b_0 = 1., sqrt(1-k**2)
  for a,b in arithmetic_geometric_mean(a_0, b_0):
    if abs(a-b) < tolerance:
      return pi/(2*a)
\end{python}

\end{frame}


\begin{frame}[fragile]{Generator expressions}
Just as we had list comprehensions, there is also \alert{generator comprehension}:
\begin{python}
g = (n for n in range(1000) if not n % 100)
# a generator that generates 0, 100, 200, 
\end{python}
\end{frame}

\begin{frame}[fragile]{Itertools}
\vfill
A useful Python module in this context is
\vskip 5pt 
\begin{python}
import itertools
\end{python}
\end{frame}



\end{document}
